\chapter{Evolución de la arquitectura}
La arquitectura presentada en el Capítulo 3 cumple su objetivo principal, demuestra que es posible controlar un dron ArduPilot desde un navegador usando Flutter. El problema aún así, son sus limitaciones estructurales. La arquitectura no llega a los requerimientos de una estación de control operativo. Este capítulo explica la evolución desde ese diseño inicial hasta la arquitectura híbrida definitiva adoptada en EZDrone, documentando los problemas técnicos que han motivado a este cambio y la implementación de sus soluciones.

\section{Limitaciones de la arquitectura base}
La arquitectura de tres capas del Capítulo 3 (Flutter Web, backend Dart y estación tierra Python) presenta tres problemas o limitaciones principales.

La primera es la latencia introducida por el polling HTTP. El frontend consulta constantemente el endpoint /status (cada 500 ms), para así obtener la telemetría y el estado del dron. Esto implica que cualquier evento puede tardar ese medio segundo en reflejarse en la interfaz. Para un sistema de monitorización a tiempo real, medio segundo puede llegar a ser aceptable, pero a la hora de controlar el movimiento del dron empieza a ser un problema serio, ya que el operador necesita una retroalimentación inmediata.

EXLICAR PROBLEMA PETICIONES HHTP

La tercera limitación, y la más estructural, es la incapacidad de soportar canales de alta frecuencia. El control por joystick, debe enviarse a 20 Hz (cada 50 ms) para poder proporcionar una respuesta relativamente fluida. Implementar esto mediante peticiones HTTP POST, generaría 1200 peticiones por minuto solamente para el joystick, cosa que saturaría tanto el backend como el broker MQTT. No solo eso sino que toda la latencia acumulada haría imposible el control del dron. El streaming de vídeo es también es directamente inviable sobre HTTP.

A estas limitaciones técnicas se añade una de diseño, el backend Dart actúa como proxy entre Flutter y el broker MQTT, añadiendo así un salto de red y un punto extra de posible fallo, y todo esto sin añadir ninguna cosa que no se pudiera implementar directamente en el propio frontend.

\section{Primera evolución: eliminación del backend Dart}
El primer paso de esta evolución es eliminar el backend Dart como intermediario y conectar Flutter directamente al broker MQTT. Esto es posible gracias al soporte de WebSocket en HiveMQ, aunque los navegadores no permiten conexiones TCP directamente, sí que permiten conexiones WebSocket.

La implementación se asienta en una clase llamada MqttLogic, directamente ubicada en el frontend. A diferencia de la arquitectura anterior que usaba MqttServerClient en el backend, esta clase inicia MqttBrowserClient con la url del broker (wss://broker.hivemp.com/mqtt) y el puerto 8884. La variante segura del WebSocket, WSS es necesaria porque la página HTTPS no permite abrir conexiones WS no seguras y el navegador directamente las bloquea.

Esta clase contiene tres operaciones principales, connect( ), que establece la conexión y registra el listener de mensajes entrantes, subscribe(topic), para suscribirse a un topic y publish(topic, payload), para publicar. Los mensajes entrantes se entregan al DronProvider, desacoplando así la capa de transporte de la lógica interna de la aplicación.

Con este cambio, el flujo de un comando de ARM (por ejemplo), pasa de cuatro saltos: Flutter - HTTP - Dart - MQTT - Python, a solamente dos: Flutter - MQTT - PYTHON. El estado de la sesión deja de depender de un proceso Dart local, y la aplicación puede desplegarse como un conjunto de archivos estáticos sin necesidad de mantener un servidor intermediario en marcha.

La gestión del estado sigue siendo el mismo modelo de estados del Capítulo 3, pero ahora el Provider es quien reacciona directamente a los topics MQTT suscritos, en lugar de tener que hacer un polling constante. Cuando la estación tierra publica, el método handleMessage del provider reacciona.

\section{Segunda evolución: Incorporación de WebRTC}
La eliminación del backend como intermediario resuelve la latencia del polling, pero no aborda nada de los canales de alta frecuencia. Para el joystick, la telemetría continua y el vídeo, se añade WebRTC, donde su arquitectura peer-to-peer, elimina el broker como figura intermediaria y proporciona la latencia y el ancho de banda necesario para estos canales.

La arquitectura de comunicación definitiva asigna a cada protocolo los canales para los que es más adecuado, basándose en las características discutidas en el Capítulo 2.

MQTT sobre WebSocket se mantiene para todos los comandos discretos (connect, arm, takeoff, land, rtl, disconnect, setMode, setCamera, detectionMode, uploadMission, startMission) y para las confirmaciones de cambio de estado. Estos mensajes son poco frecuentes, y su pérdida sería crítica.

WebRTC se usa para los tres canales mencionados anteriormente. Un DataChannel para el joystick, donde se envían los valores de los cuatro ejes desde Flutter a la estación tierra. La baja latencia de WebRTC es imprescindible para esto, ya que un retardo entre el movimiento del joystick y la respuesta del dron podría provocar problemas de control.

Otro DataChannel para la telemetría, que envía la telemetría desde la estación tierra hacia Flutter. 

El tercero es un video track que transmite el stream de la cámara del dron desde Python a Flutter. El vídeo es el caso más evidente de uso de WebRTC. No tiene alternativa razonable ya que requiere una actualización demasiado rápida para llevarse sobre MQTT o HTTP.

\section{Servidor de señalización}
La conexión WebRTC requiere un servidor de señalización para el intercambio inicial de SDP y candidatos ICE entre los dos peers. Este servidor se implementa en el backend, que queda desplegado en el propio servidor de la eetac: dronseetac.upc.edu.

El servidor es una aplicación Dart pura que levanta un servidor HTTPS sobre el puerto 8105, cargando los certificados TLS de Let’s Encrypt. Atiende dos tipos de peticiones, las que llegan a la ruta /ws se tratan como conexiones WebSocket de señalización. El resto sirven como ficheros estáticos a la carpeta /web, que contiene el build de Flutter Web compilado. Esta arquitectura permite que el mismo proceso sirva tanto la aplicación al navegador como el canal de señalización, sin necesidad de un servidor web separado.

El protocolo de señalización implementado tiene tres fases. En la primera, el cliente envía su identificador de peer “browser” para Flutter y “Python” para la estación tierra. El servidor registra el par y responde inmediatamente con la configuración ICE, que incluye un servidor STUN público y el servidor TURN privado (alojado en dronseetac.upc.edu:3478). 

En la segunda fase, los mensajes posteriores (SDP offer/answer y candidatos ICE) se reenvían al peer indicado. En la última, cuando un peer se desconecta, se elimina su entrada.

Explicar que es un servidor TURN????

\subsection{Certificación TLS}
Para que el servidor Dart pueda atender conexiones HTTPS y WSS, es necesario que disponga de un certificado TLS válido y reconocido por los navegadores. En este proyecto, el certificado se obtiene mediante Let's Encrypt, una Autoridad Certificadora gratuita y de confianza.

El proceso de emisión del certificado utiliza la herramienta certbot, que se ejecuta directamente en el servidor (en este caso en dronseetac.upc.edu). El flujo de certificación sigue los pasos siguientes:

certbot solicita un certificado a Let’s Encrypt para el dominio dronseetac.upc.edu. Let’s Encrypt emite un challenge de tipo HTTP-01, donde exige que el servidor sirva un fichero de verificación en la ruta /.well-known/acme-challenge/<token>. Certbot despliega ese fichero y notifica a Let’s Encrypt para que lo verifique. Una vez superada la verificación, Let’s Encrypt firma y emite el certificado.

El resultado del proceso son dos ficheros que certbot almacena en el sistema, fullchain.pem, la cadena completa de certificados del servidor y los certificados intermedios de la CA. Y privkey.pem, la clave privada del servidor.

Los certificados de Let’s Encrypt tienen una validez de 90 días. Para evitar interrupciones, certbot instala automáticamente una tarea que intenta la renovación cuando el certificado está a punto de caducar. Tras la renovación, los nuevos certificados reemplazan los anteriores y el servidor se reinicia para cargar los nuevos certificados. 

\section{Capa WebRTC en Flutter}
La lógica WebRTC del cliente Flutter, se encuentra dentro de la clase WebRTCLogic, en el frontend. Esta clase gestiona el ciclo de vida de la conexión, es decir, el establecimiento de la conexión, recepción de stream de vídeo, apertura de los DataChannels y desconexión.

Al llamar a connect(signalingUrl), la clase abre una conexión WebSocket hacia el servidor de señalización y envía su identificador (“browser”) como primer mensaje. A continuación, espera a la configuración ICE. Con la configuración ICE recibida, crea la RTCPeerConnection.

\section{Arquitectura híbrida definitiva}
La arquitectura resultante combina los dos protocolos de forma que cada uno opera exclusivamente en el dominio para el que es más adecuado. La arquitectura final tiene la siguiente estructura:

AÑDIR DIAGRAMA

Esta separación tiene muchas ventajas importantes. El broker MQTT solo recibe tráfico de baja frecuencia, joysticks y telemetría fluyen directamente entre los peers una vez se ha establecido la conexión WebRTC, sin tener que pasar por ningún servidor intermediario.

Una implicación operativa importante es el orden de inicio, donde la estación tierra debe iniciar la conexión MQTT al broker antes de poder recibir el comando connect, y debe de estar el canal WebRTC abierto antes de poder recibir datos o enviar telemetría. 



